\section{Conclusion}
\label{sec:conclusion}

HSAP is a falsifiable agent-based model showing that endocrine-linked behavioral feedbacks can generate population stabilization, crowding pathology, recovery, or partial collapse under reduced external threat. The model's most important finding is negative: population trajectories alone are insufficient to identify the mechanistic pathways driving these dynamics. Seven of 11 null models produced indistinguishable population outcomes. HSAP becomes testable only when population outcomes are paired with behavioral and endocrine observables.

This result has practical implications for conservation biology. Post-predator-removal monitoring programs that measure only population size may miss endocrine and behavioral signatures of regulatory change. HSAP argues for integrating endocrine sampling into population monitoring, particularly in systems where predator removal has produced unexpected population dynamics.

The model is not empirically validated. It shows that endocrine-linked behavioral feedbacks are sufficient to produce qualitatively distinct population dynamics under identical resource conditions. Whether these feedbacks are necessary, or even present, in real populations remains an empirical question.
